\documentclass[12pt,a4paper]{article}
\usepackage[utf8]{inputenc}
\usepackage[T1]{fontenc}
\usepackage{lmodern}
\usepackage[french]{babel}
\usepackage[a4paper,margin=2.3cm]{geometry}
\usepackage{amsmath,amssymb}
\usepackage{enumitem}
\setlist[itemize]{label=\textbullet,nosep}
\usepackage[most]{tcolorbox}
\usepackage{fancyhdr}
\pagestyle{fancy}
\fancyhf{}
\fancyhead[L]{\small Exercice de maison}
\fancyhead[R]{\small Suites géométriques}
\fancyfoot[C]{\small\thepage}
\renewcommand{\headrulewidth}{0.4pt}
\setlength{\parindent}{0pt}
\setlength{\parskip}{5pt}
\newcommand{\question}[1]{\par\medskip{\large\bfseries #1}\par\smallskip}

\begin{document}

\begin{center}
{\LARGE\bfseries Exercice de maison}\\[4pt]
{\large Suite géométrique}
\end{center}

\vspace{4pt}
Nom et prénom : \dotfill \quad Classe : \dotfill

\vspace{6pt}
\begin{tcolorbox}[colback=white,colframe=black,boxrule=0.5pt,arc=1pt,title=\textbf{Énoncé},
  coltitle=black,colbacktitle=gray!15]
On donne la suite $(V_n)$ définie par :
\[
\begin{cases}
V_0 = 5\\
V_{n+1} = 2V_n
\end{cases}
\]
\begin{enumerate}[nosep]
  \item Calculer $V_1$ et $V_2$.
  \item Montrer que $(V_n)$ est une suite géométrique dont on précisera la raison et le premier terme.
  \item Exprimer $V_n$ en fonction de $n$.
  \item Donner le sens de variation de $(V_n)$.
  \item Calculer $S = V_0 + V_1 + \dots + V_9$.
\end{enumerate}
\end{tcolorbox}

\end{document}
